% !TeX program = xelatex
\documentclass[12pt,a4paper]{article}
\usepackage{fontspec}
\setmainfont{Times New Roman}
\usepackage[english,russian]{babel}
\usepackage{amsmath,amssymb,mathrsfs,cancel}
\usepackage[margin=22mm]{geometry}
\usepackage{graphicx,float}
\usepackage[unicode,hidelinks]{hyperref}
\DeclareMathOperator{\grad}{grad}
\DeclareMathOperator{\Div}{div}
\newcommand{\Rey}{\mathrm{Re}}
\setlength{\emergencystretch}{3em}
\interfootnotelinepenalty=10000
\newcommand{\topic}[1]{\par\medskip\noindent\textbf{#1}\par\nobreak\smallskip}
\renewcommand{\thefootnote}{\arabic{footnote})}
\begin{document}
\section*{Вопрос 8}
\noindent Упругий режим фильтрации. Уравнение пьезопроводности. Постановка задачи о распространении плоской волны повышения давления в полубесконечном пласте. Оценка характерного времени распространения возмущения давления.

\medskip
\noindent\textbf{Источник:} Н. Е. Леонтьев. \textit{Основы теории фильтрации}. 2-е изд. Москва: МАКС Пресс, 2017. Глава III, § 12, печатные страницы 70--77.
\par\bigskip
\begin{quote}\sloppy
В отличие от газов, которые могут в разы менять свой объем, капельные жидкости (вода, нефть) являются \textit{слабосжимаемыми} веществами, у которых в промысловых условиях относительное изменение объема составляет доли процента. Поэтому, во-первых, при изучении эффектов сжимаемости на фильтрацию жидкостей нужно одновременно учитывать и сжимаемость пористого скелета и, во-вторых, учет этих изменений можно проводить в линейном приближении.

В соответствии с этим к системе, состоящей из уравнения неразрывности и закона Дарси:
\[
(m\rho)'_t+\Div(\rho\boldsymbol u)=0,\qquad \boldsymbol u=-\frac{k}{\mu}\grad p,
\]
добавляются \textit{линейные} соотношения, определяющие изотермические зависимости плотности $\rho$, пористости $m$, проницаемости $k$ и вязкости $\mu$ от возмущения давления жидкости в порах по сравнению с некоторым начальным состоянием $p_0$\footnote{Тем самым мы предполагаем деформации жидкости и пористого скелета \textit{обратимыми} (упругими), в связи с чем о процессах, при которых эта теория справедлива, говорят как об \textit{упругом режиме фильтрации}.}:
\begin{equation}\begin{aligned}
\rho(p)&=\rho_0\cdot\left(1+\frac{p-p_0}{K_\rho}\right),&m(p)&=m_0\cdot\left(1+\frac{p-p_0}{K_m}\right),\\
k(p)&=k_0\cdot\left(1+\frac{p-p_0}{K_k}\right),&\mu(p)&=\mu_0\cdot\left(1+\frac{p-p_0}{K_\mu}\right),
\end{aligned}\tag{12.1}\label{q8:eq-12-1}\end{equation}
где величины $\rho_0,m_0,k_0,\mu_0$ относятся к невозмущенному состоянию, а константы $K_\rho,K_m,K_k,K_\mu$ являются упругими модулями\footnote{Например, при $10^\circ$C для воды $K_\rho\approx2\cdot10^4$ атм, $K_\mu\approx-2\cdot10^4$ атм; отрицательное значение $K_\mu$ соответствует падению вязкости воды при не очень большом повышении давления.}, которые характеризуют изменение свойств пористой среды и жидкости при всестороннем сжатии и по порядку величины много больше характерных возмущений давления: $(p-p_0)/K\ll1$.

\setcounter{footnote}{0}
\par\nobreak\smallskip{\small Н. Е. Леонтьев. \textit{Основы теории фильтрации}. 2-е изд., 2017. Глава III, § 12, с. 70.}\par
\end{quote}

\begin{quote}\sloppy
Как и в случае фильтрации газа (§ 11), мы подставим закон Дарси в уравнение неразрывности:
\[
\frac{\partial(m(p)\rho(p))}{\partial t}=\Div\left(\frac{\rho(p)k(p)}{\mu(p)}\grad p\right),
\]
что после преобразований дает
\[
\rho m\left(\frac{m'_p}{m}+\frac{\rho'_p}{\rho}\right)p'_t=\frac{\rho k}{\mu}\left\{\left(\frac{\rho'_p}{\rho}+\frac{k'_p}{k}+\frac{(1/\mu)'_p}{1/\mu}\right)|\grad p|^2+\Delta p\right\},
\]
и, поскольку мы хотим получить уравнения для \textit{малых возмущений}, упростим коэффициенты при производных от давления: вместо значений $k,\mu,\rho$ и $m$ возьмем их невозмущенные значения $k_0,\mu_0,\rho_0$ и $m_0$, а при вычислении производных воспользуемся линейными уравнениями состояния (12.1), например:
\[
\frac{\rho'_p}{\rho}\approx\frac{1}{K_\rho},\qquad \frac{(1/\mu)'_p}{1/\mu}\approx-\frac{1}{K_\mu}.
\]
В результате получится уравнение с постоянными коэффициентами, в котором дополнительно можно отбросить нелинейное слагаемое\footnote{Зачеркнутый член имеет порядок $(p-p_0)^2/(KL^2)$, где $L$ --- характерный размер задачи, а оператор Лапласа порядка $(p-p_0)/L^2$, так что отношение этих членов порядка $(p-p_0)/K\ll1$.}:
\[
\rho_0m_0\left(\frac{1}{K_m}+\frac{1}{K_\rho}\right)p'_t=\frac{\rho_0k_0}{\mu_0}\left\{\cancel{\left(\frac{1}{K_\rho}+\frac{1}{K_k}-\frac{1}{K_\mu}\right)|\grad p|^2}+\Delta p\right\},
\]
что в итоге дает \textit{уравнение пьезопроводности}\footnote{Название (от греч. $\pi\iota\varepsilon\zeta\omega$ --- «давлю») дано по аналогии с названием уравнения теплопроводности, совпадающего по форме с (12.2).}, полученное рядом исследователей в 1930--40-х гг. [11,17]:
\begin{equation}
p'_t=\varkappa\Delta p,\qquad \varkappa=\frac{k_0}{m_0\mu_0}\left(\frac{1}{K_\rho}+\frac{1}{K_m}\right)^{-1}=\mathrm{const},\tag{12.2}\label{q8:eq-12-2}
\end{equation}
где $\varkappa$ --- \textit{коэффициент пьезопроводности}. Граничные условия для уравнения (12.2) ставятся аналогично условиям при изотермической фильтрации газа (стр. 67).

\setcounter{footnote}{0}
Коэффициент пьезопроводности одновременно учитывает сжимаемость жидкости и пористого скелета. Отметим, что в промысловых условиях встречаются ситуации, когда фильтрующаяся жидкость содержит пузырьки газа (например, газ, растворенный в жидкости, может выделяться при понижении давления). Это приводит к увеличению средней сжимаемости флюида и, как следствие, к резкому снижению эффективного коэффициента пьезопроводности $\varkappa$.
\par\nobreak\smallskip{\small Н. Е. Леонтьев. \textit{Основы теории фильтрации}. 2-е изд., 2017. Глава III, § 12, с. 71--72.}\par
\end{quote}

\begin{quote}\sloppy
Рассмотрим одномерную задачу движении плоской волны давления в полубесконечном пористом пласте. Предположим, что начальное давление флюида в пласте постоянно, а в нулевой момент времени давление на границе пласта \textit{скачком} повышается на постоянную величину $\delta p$ и далее все время поддерживается на этом уровне. Считая для простоты начальное давление нулевым: $p_0=0$\footnote{В уравнение пьезопроводности входят только производные от давления.} и направляя ось $Ox$ вдоль пласта, получим задачу
\begin{equation}
p'_t=\varkappa p''_{xx},\quad p(x,0)=0,\quad p(0,t)=\delta p,\quad \lim_{x\to\infty}p(x,t)=0\tag{12.3}\label{q8:eq-12-3}
\end{equation}
(начало координат выбрано на границе пласта). Решение задачи автомодельно и, в соответствии с $\pi$-теоремой, имеет вид
\[
p=p(x,t,\varkappa,\delta p)=\delta p\cdot f(\xi),\qquad \xi=\frac{x}{\sqrt{\varkappa t}},
\]
что сводит (12.3) к краевой задаче для обыкновенного дифференциального уравнения
\[
2f''(\xi)+\xi f'(\xi)=0,\quad f(0)=1,\quad f(\infty)=0,
\]
решение которой выражается через функцию ошибок (рис. 27 \textit{а}):
\[
f(\xi)=1-\operatorname{erf}(\xi/2),\qquad \operatorname{erf}x=\frac{2}{\sqrt\pi}\int_0^x\exp(-\alpha^2)\,d\alpha.
\]
В размерных переменных распределение давления дается выражением
\[
p(x,t)=\delta p\cdot\left(1-\operatorname{erf}\left(\frac{x}{2\sqrt{\varkappa t}}\right)\right);
\]
качественный вид пространственного распределения давления для трех моментов времени $t_1<t_2<t_3$ показан на рис. 27 \textit{б}.


Скорость фильтрации, которая находится с помощью закона Дарси:
\[\begin{aligned}
u_x(x,t)&=-\frac{k(p(x,t))}{\mu(p(x,t))}p'_x(x,t)\approx-\frac{k_0}{\mu_0}p'_x(x,t)=\\
&=\frac{k_0\delta p}{\mu_0}\frac{1}{\sqrt{\pi\varkappa t}}\exp\left(-\frac{x^2}{4\varkappa t}\right),
\end{aligned}\]
в каждой фиксированной внутренней точке пласта с течением времени сначала возрастает от нулевого значения до некоторого максимального, а затем постепенно падает до нуля при $t\to\infty$ (рис. 27 \textit{в}).

С помощью этого решения несложно оценить характерное время $\tau$ переходного процесса в упругом режиме: $\tau\sim L^2/\varkappa$. При характерных параметрах $K\sim10^4$ атм, $k\sim10^{-8}$ см$^2$, $\mu\sim10^{-2}$ г/(см$\cdot$с), $m\sim1$ получим $\varkappa\sim1$ м$^2$/с, и если $L\sim10$ м, то $\tau\sim1$ мин, а при $L\sim10^3$ м $\tau\sim10$ сут.
\par\nobreak\smallskip{\small Н. Е. Леонтьев. \textit{Основы теории фильтрации}. 2-е изд., 2017. Глава III, § 12, с. 72--73.}\par
\end{quote}

\end{document}
